% =====================================================================
%  Capítulo 9 — Limitações e próximo ciclo
% =====================================================================
\chapter{Limitações e próximo ciclo}
\label{cap:conclusao}
\rubricaComunicacao

\begin{orientacao}
Uma página. Sem repetir o resumo e sem adjetivo de autoelogio.
\end{orientacao}

\section{Limitações conhecidas}

\begin{orientacao}
O que o sistema não faz, o que só funciona sob condições controladas e
qual número do relatório é o mais frágil. Todo projeto deste ciclo tem
pelo menos um número que é estimativa: identifique-o. As limitações
abaixo vêm da AP1; mantenham só as que continuam valendo e acrescentem
as da AP2.
\end{orientacao}

\textbf{Herdadas da AP1, e que continuam valendo.} A autonomia da
\RNF{01} continua calculada, e não medida: os 21 dias de agora usam o
tempo de rádio medido, mas nenhum ensaio com bateria real foi feito, e o
consumo do alarme segue fora da conta. A robustez a interferência
eletromagnética (\RNF{05}) ganhou só uma evidência indireta, os 33
minutos de ar-condicionado ligado na campanha, sem o ensaio dedicado
junto a um compressor. A lógica de mudança de faixa (\RF{02}) continua
falhando nos testes de bancada da AP1, a cessação do alarme (\RF{04})
continua sem log formal, e o sensor desconectado ainda leva o nó a um
laço de reinícios pelo \emph{watchdog} (\CT{05}). Nenhuma delas foi
tratada neste ciclo.

\textbf{Da AP2.} A campanha durou \SI{2,78}{\hour}, numa tarde de
24/10, num cômodo com ar-condicionado e não no cenário de referência, com
um nó na USB e o roteador no mesmo ambiente; ela não mostra o ciclo
diário, a madrugada, as paredes da farmácia nem a porta do freezer. As
demais limitações conhecidas:

\begin{itemize}
  \item \textbf{O nó que publica não é o nó que alarma.} O
    \emph{firmware} da campanha (branch \texttt{lab3}) não tem a leitura
    filtrada, o alarme local, a configuração na NVS nem o \emph{deep
    sleep} da AP1; o \emph{firmware} da AP1 não publica. Tudo o que o
    Cap.~\ref{cap:testes} mede de comunicação vale para o primeiro, e
    tudo o que mede de alarme, para o segundo.
  \item \textbf{Sem fila local e sem \texttt{boot}.} Toda queda de
    conexão vira leitura perdida de vez (cerca de 740 na campanha), e um
    reinício só é reconhecido se a primeira leitura depois dele chegar.
  \item \textbf{A causa da fase degradada é uma hipótese.} O
    \emph{timeout} de \SI{1}{\second} da biblioteca MQTT explica o padrão
    das quedas, mas a latência até o \emph{broker} não foi medida e o
    teste com \emph{timeout} maior não foi feito. O número mais frágil
    da AP2, depois da corrente do rádio, é a própria completude: outra
    tarde, com outra carga na rede, poderia dar 96\,\% ou menos que
    63\,\%.
  \item \textbf{\emph{Broker} público, sem autenticação.} Qualquer
    pessoa poderia ter publicado em \texttt{ufpi-iot26/eq01/\#}; nada
    estranho chegou, mas nada impedia.
  \item \textbf{Período de \SI{5}{\second}, e não de \SI{600}{\second}.}
    A campanha não exercita o período do \RF{01} nem mede a energia do
    nó sempre conectado.
  \item \textbf{O número mais frágil é a corrente do rádio.} Os
    \SI{300}{\milli\ampere} da Seção~\ref{sec:custo-radio} são da placa
    inteira, lidos no medidor USB, e não do ESP32; com a estimativa
    anterior de \SI{120}{\milli\ampere}, a autonomia calculada passaria de
    21 para 26 dias. O tempo de rádio, ao contrário, foi medido com
    \texttt{millis()} em 36 ciclos.
  \item \textbf{A conexão TLS foi medida contra um servidor de eco, não
    contra o \emph{broker} da campanha}, e a corrente do nó conectado e
    ocioso, que é o que pesa num nó que não dorme (\DP{6}), não foi
    medida.
  \item \textbf{O alcance foi medido num ambiente interno de teste, com
    um único teste de \SI{60}{\second} por posição}, sem caracterizar as
    paredes nem controlar a orientação das placas; a contagem de paredes
    digitada a 8 e \SI{11}{\metre} difere entre as baterias de ESP-NOW e
    de BLE, embora o RSSI mostre que as posições foram as mesmas
    (Apêndice~\ref{ap:testes}, \CT{15}).
    Câmaras frias e freezers, cujas paredes metálicas devem atenuar muito
    mais, não foram testados.
  \item \textbf{Os testes de robustez da AP1 (\CT{05} e \CT{06}) foram
    feitos com um DHT11}, e não com o DHT22 do projeto; o comportamento
    observado (laço de reinícios e recuperação da NVS) não depende do
    sensor, mas a precisão declarada nas mensagens, sim.
  \item \textbf{O nó não confere o certificado do \emph{broker}}
    (\texttt{setInsecure()}, herdado do Lab~03): a conexão é cifrada, mas
    não autenticada.
  \item \textbf{LoRa/LoRaWAN, a tecnologia escolhida para o produto, não
    foi testado}: o módulo comprado pela equipe não chegou a tempo. O
    alcance, o tempo no ar e o consumo do LoRa vêm da matriz de 22/09, da
    folha de dados e da conta do ar da Aula~14, e não de medição; a
    autonomia de $\approx$29 dias do nó de borda é estimativa. Os
    substitutos testados no lugar dele foram medidos numa casa, e não na
    farmácia: o ESP-NOW foi bem até \SI{8}{\metre} com uma parede, mas
    pode não repetir esse resultado no nosso caso de uso, e o interior do
    freezer não foi testado com nenhum rádio. O nó central continua sem hardware
    definido.
\end{itemize}

\section{O que entra no próximo ciclo}

\begin{orientacao}
Ligue este relatório ao seguinte.
\aPartirDaAP{1}{}%
\seAP{1}{Na AP2 a equipe vai integrar este sistema a uma infraestrutura
de comunicação e formar o conjunto de dados do projeto: diga qual
tecnologia é a candidata e qual linha do orçamento de energia vai ser
medida em vez de estimada.}%
\seAP{2}{Na AP3 a equipe vai treinar e embarcar um modelo: diga o que
pretende classificar ou prever, e se o conjunto de dados formado neste
ciclo é suficiente para isso.}%
\seAP{3}{Diga o que o projeto precisaria para sair do laboratório e
virar algo instalado no contexto do Capítulo~\ref{cap:problema}:
alimentação, invólucro, manutenção, custo por unidade em escala.}
\end{orientacao}

Na AP3 a equipe vai treinar e embarcar um modelo sobre o conjunto de
dados formado neste ciclo. O que se pretende é detectar uma falha de
climatização --- o ar-condicionado que desliga, ou a porta de um
freezer que fica aberta --- pela forma da curva de temperatura e
umidade, antes que a leitura saia da faixa. O conjunto tem \textbf{um} exemplo desse evento, anotado no
registro do Apêndice~\ref{ap:campanha} (16:35 de 24/10), o que não
basta para treinar nem para validar. Falta coletar mais
desligamentos, em dias e horários diferentes e de preferência no
cenário de referência, e trechos longos sem evento.

Antes disso, o \emph{firmware} precisa de três mudanças, nesta ordem:
uma fila local, para que uma queda de conexão atrase a leitura em vez
de apagá-la; o \texttt{boot} no contrato; e a união com o
\emph{firmware} da AP1, para que o nó que publica seja o mesmo que
alarma. Junto com elas, o teste do \emph{timeout} da biblioteca MQTT
(\texttt{mqtt.setTimeout(5000)}) com a contagem de \texttt{offline}
antes e depois. Quando o módulo LoRa chegar, o enlace entre o nó de
borda e o nó central passa a ser LoRa, com o contrato em binário de 8
bytes, e o alcance e o consumo dele serão medidos nas mesmas posições
das baterias de ESP-NOW e BLE, de preferência também com o nó dentro de
um freezer.

A ponte com o servidor do produto (\CT{16}) resolveu duas questões do
contrato e deixou duas para o próximo ciclo. Resolvidas: a hora da
leitura continua sendo a do nó (o \texttt{ts} do NTP vira o
\texttt{data\_hora} em UTC), e o \texttt{lote\_id} derivado do nó e do
intervalo de \texttt{seq} torna o reenvio idempotente, como o reenvio dos
13 lotes mostrou. Abertas: no produto, o nó de borda em \emph{deep sleep}
com LoRa não terá NTP, e a hora terá de vir do nó central, com o
atraso declarado no contrato; e o \texttt{boot} precisa entrar no
\texttt{lote\_id}, para que um reinício não gere o mesmo identificador
para leituras diferentes. O servidor tem ainda um defeito a corrigir:
um lote com várias leituras fora da faixa em temperatura e umidade
derruba a abertura de alertas (HTTP 500), e o lote é recusado inteiro.

\section{Divisão do trabalho}

\begin{orientacao}
Uma linha por integrante, dizendo o que cada um fez neste ciclo. Os
papéis rotacionam a cada módulo: ninguém repete na AP2 o papel que teve
na AP1. Na AP1: João Pedro na integração; Victor e Davi no
\emph{firmware}; Felipe no \emph{hardware} e nos testes; Luanderson na
documentação.
\end{orientacao}

Os papéis rotacionaram em relação à AP1, e cada integrante ficou com
uma frente. Os \emph{commits} da AP2 foram enviados ao repositório por
um único integrante (João Pedro), que concentrou a inserção do material
de todos; por isso a autoria das tarefas é a desta lista, e não a do
histórico do Git.

\begin{itemize}
  \item \textbf{Felipe} (\emph{firmware} e comunicação; na AP1,
    \emph{hardware} e testes): \emph{firmware} MQTT da campanha sobre o
    sketch do Lab~03 (DHT22 no \texttt{GPIO4}, tópicos da equipe, NTP,
    testamento e sessão persistente); \emph{firmware} de sensor e de
    \emph{gateway} das baterias de ESP-NOW e de BLE; ferramenta da
    Aula~10 que cronometra o ciclo de rádio (branches \texttt{lab3},
    \texttt{teste-esp-now}, \texttt{teste-ble} e
    \texttt{apendice\_b}).
  \item \textbf{João Pedro} (dados --- gravador e qualidade; na AP1,
    integração): mudanças no \texttt{gravador.py} e no
    \texttt{qualidade.py}, \texttt{grafico.py} novo, gravadores seriais
    das baterias e reanálise dos dados brutos (\texttt{analise.py});
    hospedou os testes de falha e a campanha, com o nó, o roteador e o
    notebook do gravador, e anotou os eventos do ambiente.
  \item \textbf{Luanderson} (dados --- campanha e registro de eventos;
    na AP1, documentação): roteiro dos testes de falha e da campanha,
    ficha da campanha, diário e registro de eventos
    (Apêndice~\ref{ap:campanha}), contagens sobre o CSV (fases, quedas de
    conexão, atraso) e consolidação dos números usados no relatório.
  \item \textbf{Victor} (integração --- repositório, \emph{broker} e
    serviços de telemetria; na AP1, \emph{firmware}): organização do
    repositório e das branches de trabalho, \emph{commit} da entrega com
    a \emph{tag} \texttt{ap2}, contrato de tópicos e configuração do
    \emph{broker}, e a ligação com o servidor do produto (nó central,
    lote e chave de API).
  \item \textbf{Davi} (documentação --- relatório, registro de decisões e
    prazos; na AP1, \emph{firmware}): migração do relatório para o modelo
    da AP2, caixas de decisão do Módulo~2 (\DP{6} a \DP{8}), tabela de
    rastreabilidade, declaração de IA e controle dos prazos da entrega.
\end{itemize}

\imprimirTabelaIntegrantes
